\section{Preliminaries}\label{sec:prelim}

We start with notation and basic concepts.

\subsection{Standard neural network notation}\label{subsec:relu_notation}

For simplicity, we start with
standard feedforward networks with $L$ hidden layers without bias.
Let $d^{(0)}, \ldots, d^{(L)}$ denote the layer widths, with input dimension $d^{(0)}$ and output dimension $d^{(L)}$, and
weight matrices $W^{(l)} \in \RR^{d^{(l)} \times d^{(l-1)}}$.
One can include bias by using an \emph{intercept input} of~$1$. 

Using an activation function $\mathtt{act}$, we denote by $a^{(0)} \coloneqq x \in \RR^{d^{(0)}}$ the input layer, by $z^{(l)} \coloneqq W^{(l)} a^{(l-1)}$ the pre-activation and by $a^{(l)} \coloneqq \mathtt{act}(z^{(l)})$ the activation.
Furthermore, we denote the function represented by the layers $l, \dots, L$ of the original network by $f^{(l)} : \ \RR ^ {d^{(l)}} \longrightarrow \RR$. 

\subsection{Convex and monotone neural networks}

We call a neural network with convex and monotone activation functions \emph{input-convex neural network} (ICNN) if all weights except possibly the first layer are non-negative~\citep{Amos2017}.
We call a neural network with monotone activation functions \emph{monotone} if all its weights are non-negative. 
It is well known that indeed ICNNs represent convex functions and monotone networks represent monotone functions; see \ref{app:fundamental}.
Note that both Maxout and ReLU are monotone and convex activation functions.